\documentclass[letterpaper]{article}
\usepackage[submission]{aaai2027}
\makeatletter
\gdef\showauthors@on{T} % W3PHIAI-27 is single-blind: reviewers see author names
\makeatother
\usepackage[hyphens]{url}
\usepackage{graphicx}
\urlstyle{rm}
\def\UrlFont{\rm}
\usepackage{natbib}
\usepackage{caption}
\frenchspacing
\usepackage{booktabs}
\setcounter{secnumdepth}{0}

\pdfinfo{
/TemplateVersion (2027.1)
}

\title{COMPASS: A Retrieval-Augmented Generation System for Evidence-Grounded Opioid Use Disorder Guidance in Tennessee}
\author{
    Emma Baumgartner\textsuperscript{\rm 1},
    Tom Berg\textsuperscript{\rm 1},
    Gerald Jones Jr.\textsuperscript{\rm 1}
}
\affiliations{
    \textsuperscript{\rm 1}Applied Systems Laboratory, Department of Industrial \& Systems Engineering,\\
    University of Tennessee, Knoxville
}
\copyrighttext{Submitted to W3PHIAI-27 (Poster Track) for review. Not yet published; copyright and publication details will appear in the final version if accepted.}

\begin{document}

\maketitle

\begin{abstract}
Tennessee has one of the highest opioid overdose death rates in the country, yet the clinical guidelines, research evidence, surveillance data, and policy documents needed to act on it are scattered across disconnected sources. We propose COMPASS, a retrieval-augmented generation (RAG) system intended to ground natural-language answers about opioid use disorder (OUD) in a curated, Tennessee-focused document corpus, returning cited sources with every response. Alongside retrieval, we plan to construct a hand-authored concept map intended to add structured relational context on top of plain retrieval. This paper describes the motivation and planned system design, along with the evaluation we intend to run: a comparison of plain retrieval against concept-map-augmented generation and an evidence re-ranking method, with repeated baseline measurements planned before any comparison is interpreted, and a demographic-equity pilot to test whether clinically irrelevant patient wording affects retrieval or answer quality. No results are reported here; this is a description of planned work.
\end{abstract}

\section{Motivation}

Tennessee has one of the highest opioid overdose death rates in the country \citep{tdh2025}. The information needed to treat, prevent, or plan around opioid use disorder (OUD) does not live in one place: clinical practice guidelines, the research evidence base for medications for OUD, state-level surveillance data, and policy documents are each held by different sources with no unified, citable point of access. A clinician or public health planner who wants a grounded answer to a concrete question currently has to search these sources independently, with no way to know if they found everything relevant.

\section{Related Work}

Retrieval-augmented generation grounds LLM generation in retrieved passages, reducing hallucination and enabling source attribution \citep{lewis2020rag}; COMPASS is planned around this now-standard retrieve-then-read pattern. Combining knowledge graphs with RAG has been explored to provide multi-hop reasoning and relationship context \citep{edge2024graphrag}, and a near-identical architecture, RAG plus an LLM plus a concept map, was independently built for Scotland's medication-assisted treatment standards \citep{scotlandmat2026}, which explicitly lists a graph-vs-plain-RAG ablation as future work; our planned evaluation is designed to run exactly that comparison. \citet{clinicbot2026} argue that clinical RAG systems should rank evidence by clinical significance rather than text similarity alone, motivating the evidence re-ranker we plan to test alongside the concept map. On the equity side, the literature is mixed: a fairness audit of machine-learning models for OUD treatment retention found real subgroup bias \citep{wang2026fairness}, while a study of GPT-4 and Gemini found no demographic effect on opioid-recommendation decisions \citep{young2024}. This disagreement is itself the motivation for testing the question directly on COMPASS rather than assuming either result generalizes.

\section{Proposed System}

We propose to address this with COMPASS, a retrieve-then-read RAG pipeline: a curated, screened document corpus spanning clinical, behavioral, and policy topic areas with particular depth on Tennessee, embedded and indexed for dense retrieval, paired with a large language model that generates answers only from retrieved passages and cites its sources. Alongside the retrieval pipeline, we plan to build a hand-authored concept map encoding relationships among medications, conditions, guidelines, and Tennessee-specific policy, intended to supply the system with relational context that plain chunk-level retrieval cannot represent on its own, such as which medications are contraindicated for which conditions.

\section{Planned Evaluation}

We plan to evaluate the system with an independent LLM judge scoring faithfulness, answer relevance, and completeness over a held-out test set of representative questions. Our central planned comparison is whether injecting concept-map context into generation, and whether a ClinicBot-inspired evidence re-ranker \citep{clinicbot2026} that reorders retrieved passages by content type rather than similarity alone, measurably improve answer quality over plain retrieval. Because small-scale RAG evaluations can show run-to-run variation large enough to make a single comparison misleading, we plan to repeat the unmodified baseline configuration several times to establish a noise floor before interpreting any difference between conditions as real.

We also plan a demographic-equity pilot: a set of matched clinical vignettes, each posed as a base patient plus variants that change only demographic wording (race, gender, insurance, county, criminal-justice involvement) while holding the clinical question identical, to test whether such wording measurably changes either the documents retrieved or the quality of the generated answer. Should a retrieval-stage sensitivity to irrelevant wording be found, we plan to test a deterministic mitigation that removes such wording before retrieval while preserving it for generation, and to verify with a separate noise-floor measurement whether any residual difference reflects a real effect or ordinary generation stochasticity.

\section{Anticipated Limitations}

We expect a modestly sized test set and a single LLM judge to limit how small an effect this evaluation can reliably detect, and we plan to size our claims accordingly rather than over-interpret single-run results. We anticipate that a hand-authored concept map built before any automated expansion will be small relative to the space of clinically relevant relationships, and plan to treat early results from it as provisional. No clinician review is planned in this first pass; we see that as an important follow-on step once an automated evaluation has identified which comparisons are worth a domain expert's time. We present this proposal to get early feedback on the evaluation design before committing further resources to it.

\small
\bibliographystyle{aaai2027}
\begin{thebibliography}{99}

\bibitem[{Edge et al., 2024}]{edge2024graphrag} Edge, D., et al. 2024. From local to global: A graph RAG approach to query-focused summarization. \emph{arXiv:2404.16130}.

\bibitem[{Lewis et al., 2020}]{lewis2020rag} Lewis, P., et al. 2020. Retrieval-augmented generation for knowledge-intensive NLP tasks. \emph{NeurIPS}, 33.

\bibitem[{Nananukul and Kejriwal, 2026}]{clinicbot2026} Nananukul, N.; and Kejriwal, M. 2026. ClinicBot: A guideline-grounded clinical chatbot with prioritized evidence RAG and verifiable citations. \emph{arXiv:2605.00846}.

\bibitem[{Nwobi et al., 2026}]{scotlandmat2026} Nwobi, S. C.; Loukil, Z.; and Jawahar, A. 2026. Development of an AI-driven chatbot for medication-assisted treatment standards in Scotland. \emph{Frontiers in Digital Health}, 8: 1698657.

\bibitem[{Tennessee Department of Health, 2025}]{tdh2025} Tennessee Department of Health, Office of Informatics and Analytics. 2025. \emph{2023 Tennessee Drug Overdose Deaths}.

\bibitem[{Wang et al., 2026}]{wang2026fairness} Wang, T.; Vivas-Valencia, C.; Bauer, C.; Gong, Y.; Choo, K.-K. R.; and Guo, Y. 2026. Toward fairness in machine learning models for predicting treatment retention and premature discontinuation in medication for opioid use disorder. \emph{arXiv:2609.22113}.

\bibitem[{Young et al., 2024}]{young2024} Young, C. C.; Enichen, E.; Rao, A.; and Succi, M. D. 2024. Racial, ethnic, and sex bias in large language model opioid recommendations for pain management. \emph{Pain}, 166(3): 511--517.

\end{thebibliography}

\end{document}
